% MANUAL RAPIDO: LABEL, REF, EQREF E HYPERREF
%%%%%%%%%%%%%%%%%%%%%%%%%%%%%%%%%%%%%%%%%%%%%%

%================================================
% 1. \label{nome}
%=================================================
% Cria um marcador interno no LaTeX.
% Esse marcador serve para guardar o numero de uma secao, equacao,
% teorema ou outro elemento numerado.
%-------------------------------------------------
% Exemplo:
%
%       \section{Introducao}
%       \label{sec:intro}
%
% Aqui, o nome sec:intro fica associado ao numero da secao.
%
%=======================================================
% 2. \ref{nome}
%=======================================================
% Recupera o numero salvo pelo \label.
% Ele mostra apenas o numero, sem parenteses.
%------------------------------------------------------
% Exemplo:

%       Veja a Seção \ref{sec:intro}.
%
% Se a secao acima for a secao 1, o resultado sera:
% Veja a Seção 1.
%
%=======================================================
% 3. \eqref{nome}
%=======================================================
% Recupera o numero de uma equacao com parenteses.
% Esse comando e usado para equacoes numeradas.
% Para funcionar, e preciso carregar o pacote amsmath.
%-------------------------------------------------------
% Exemplo:
%       \begin{equation}
%       \label{eq:principal}
%       x^2+y^2=1
%       \end{equation}
%
%       Veja a equação \eqref{eq:principal}.
%
% Se essa for a equacao 3, o resultado sera:
% Veja a equação (3).
%
%=======================================================
% 4. \usepackage{hyperref}
%=======================================================
% Esse pacote faz as referencias virarem links clicaveis no PDF.
% Ou seja, ao clicar em \ref ou \eqref no PDF,
% o leitor vai direto para a secao ou equacao correspondente.
%---------------------------------------------------------
% Exemplo no preambulo:
%
%       \usepackage{hyperref}
%
%
% 5. EXEMPLO COMPLETO
%
%       \usepackage{amsmath}
%       \usepackage{hyperref}
%
%       \section{Exemplo}
%       \label{sec:exemplo}
%           |
%           |    \begin{equation}
%           |    \label{eq:teste}
%           |     |   x^2+y^2=1
%           |   \end{equation}
%           |
%           |   Veja a Seção \ref{sec:exemplo}.
%           |   Veja a equação \eqref{eq:teste}.
%